\chapter{Előfeldolgozás}

Az előfeldolgozás fejezetben helyet kap az adathalmaz leírása, tervezés és implementálás bemutatása az adathalmaz letöltésétől egészen a tokenizált adatok mentéséig.
A tervezés a végleges változat dokumentálását fedi le, és jellegzetesebb kihívások és megoldások csak említés szintjén kerülnek megörökítésre. 
Az implementálás ötletesebb megoldások bemutatására szorítkozik, és rövid leírására az alkalmazott módszereknek.

\section{Adathalmaz eredete}

A felhasznált adathalmaz M. Wiechmann \texttt{enron\_spam\_data} GitHub-oldalán található vesszővel tagolt adatfájl 
(\texttt{CSV}) formátumban~\cite{wiechmannNoppelEnronSpam}.

Az eredeti adathalmazról, mely abban különbözik a leírtak szerint, hogy minden egyes e-mail külön fájlként van tárolva, 
a Spam filtering with Naive Bayes – Which Naive Bayes? tanulmányban olvashatunk a harmadik fejezetben. 
A következő alfejezet ezt a művet veszi alapul az adathalmaz általános bemutatására~\cite{inproceedings}.

Az Enron botrány kivizsgálásának keretében közel 150 Enron dolgozó fájljai lettek közzé téve, 
melyek között jelentős mennyiségű személyes e-mail üzenet is megtalálható. 
Személyre szabott spam filtert tűztek ki célul, így kiválasztottak 6 Enron dolgozót: 
farmer-d, kaminski-v, kirchen-l, williams-w3, beck-s és lokay-m, a Bekkerman által szolgáltatott adathalmazból, 
melyben a levelezések csak ham üzeneteket tartalmaztak ~\cite{inproceedings}.

A Spam üzeneteket négy különböző forrásból állították össze melyek sorra: a SpamAssassin corpus, a Honeypot project, 
Bruce Guenter spam gyűjteménye, és a 3. fejezet bevezetésében említett mű harmadik szerzőjének, Paliouras G. egyéni spam gyűjteménye ~\cite{inproceedings}.

A hat darab ham üzenet gyűjteményből mindegyik párba lett állítva egy spam gyűjteménnyel a három közül, 
és az arányuk három esetében 1:3-hoz lett, míg a másik három esetében 3:1-hez ~\cite{inproceedings}.

A következő előfeldolgozási műveleteket is elvégezték:
\begin{itemize}
\item Az e-mail fiók tulajdonosa által küldött levelek eltávolítása
\item HTML címkék eltávolítása
\item E-mail fejléc eltávolítása
\item Nem-latin karaktereket alkalmazó spam-ek eltávolítása
\end{itemize}

\section{Előfeldolgozás tervezése}

Egy adathalmaz előfeldolgozása során több iterációban történt az információ kinyerése az adatokról, 
majd ezen információk alapján előfeldolgozási módosítások elvégzése. A következő folyamatábra és 
leírás dokumentálja a folyamatot.